%!TEX root = ./../../springer-mono-format/book.tex
%!TEX root = ./../../thesis-trier.tex

\IN{permutation}
\IN{inverse transformation}
A \emph{permutation} is a bijective transformation.
Recall that, in this case, as discussed in \cref{prelim:subsec:func:basic}, 
there exists an \emph{inverse transformation} $t^{-1} \colon X \to X$
such that $\concat{t}{t^{-1}} = \concat{t^{-1}}{t} = \operatorname{id}_X$.


\IN{permuted subset} 
\IN{quasipermutation}
For an arbitrary transformation $t \colon X \to X$
and $S \subseteq X$ we say that $t$ \emph{permutes a subset $S \subseteq X$} 
if $t_{|S}$ is a permutation on $S$. 
We call~$t \colon X \to X$ a \emph{quasipermutation}~\cite{GorKou95} if it permutes its image $t(X)$.

First, in~\cref{prelim:prop:permutation} we collect several well-known characterizations of permutations
on finite sets. Then, the remaining statements investigate quasipermutations.

\begin{proposition}
	\label{prelim:prop:permutation} \label{prelim:prop:non-perm-dupl-cycle}
	Let $t \colon [n] \to [n]$. Then the following are equivalent.
	\begin{enumerate}
		\item \label{prelim:prop:item:not_perm} $t$ is a permutation.
		
		\item \label{prelim:prop:item:non_injective} $\dupl{t} = \emptyset$, i.e., $t$ is injective.		
		
		\item \label{prelim:prop:item:invdupl} $t^{-1}(\dupl{t}) \cap t([n])$ is empty.
				
		\item \label{prelim:prop:item:root} $t$ has no root point.
				
		\item \label{prelim:prop:item:defect} $\defect{t} = 0$.
		
		\item \label{prelim:prop:item:cycles} Every state is a cycle state.
		
		\item \label{prelim:prop:item:identity} There exists $i > 0$ such that $t^i = \operatorname{id}_{[n]}$.			
		
		\item \label{prelim:prop:item:not_surjective} $\rank{t} = n$, i.e., $t$ is surjective.		
	\end{enumerate}	
\end{proposition}
\begin{proof} 
	$\ref{prelim:prop:item:not_perm}\Rightarrow~\ref{prelim:prop:item:non_injective}$: This follows as a permutation 
	is injective by definition.
		
	$\ref{prelim:prop:item:non_injective}\Rightarrow~\ref{prelim:prop:item:invdupl}$: 
	This is obvious.
	
	$\ref{prelim:prop:item:invdupl}\Rightarrow~\ref{prelim:prop:item:root}$:
	Let $y \in [n]$ be a root point.
	Then $y \in \dupl{t}$ and $t^p(y) = y$ for some $p > 0$.
	This gives $t^{p-1}(y) \in t^{-1}(\dupl{t}) \cap t([n])$.	
		
	$\ref{prelim:prop:item:root}\Rightarrow~\ref{prelim:prop:item:defect}$:
	Suppose $\defect{t} > 0$.
	Then there exists $x \in \excl{t}$.
	By finiteness, there exists a minimal $i \ge 0$
	such that $t^i(x) = t^{i+p}(x)$ for some $p > 0$, where, since $x \in \excl{t}$, we must have $i > 0$.
	Set $y = t^{i}(x)$. We have $t^p(y) = y$.
	By the minimal choice of~$i$, we have $t^{i-1}(x) \ne t^{i + p - 1}(x)$
	and both map to the same point~$y$. Hence we have $y \in \dupl{t}$.
	Thus~$y$ is a root point.
		
	$\ref{prelim:prop:item:defect}\Rightarrow~\ref{prelim:prop:item:cycles}$:
	By \cref{prelim:prop:item:defdupl2} from~\cref{prelim:prop:gentransdef},
	we have $|\dupl{t}| \le \defect{t} = 0$, i.e.~$t$ is injective.
	By finiteness
	there exist $i \ge 0$ and $p > 0$ with $i$ minimal
	such that $t^i(x) = t^{i+p}(x)$.
	For $i > 0$, by injectivity, we obtain $t^{i-1}(x) = t^{i - 1 + p}(x)$,
	contradicting the minimal choice of $i$. So $i = 0$ and $x = t^p(x)$. Thus~$x$
	is a cycle state.
		
	$\ref{prelim:prop:item:cycles}\Rightarrow~\ref{prelim:prop:item:identity}$: 
	By assumption for each $x \in [n]$ there exists $i_x > 0$
	such that $\applyfunc{x}{t^{i_x}} = x$.
	Then $t^{\lcm\{ i_x \mid x \in [n] \}} = \operatorname{id}_{[n]}$.	
	
	$\ref{prelim:prop:item:identity}\Rightarrow~\ref{prelim:prop:item:not_surjective}$: 
	By assumption we have $x = t(t^{i-1}(x)) \in t([n])$ for each $x \in [n]$.
	So $t([n]) = [n]$, which gives $\rank{t} = n$.	
	
	$\ref{prelim:prop:item:not_surjective}\Rightarrow~\ref{prelim:prop:item:not_perm}$:
	By \cref{prelim:prop:item:rank-defect-eqn} from \cref{prelim:prop:rank-defect} 
	we have $n = \rank{t} = n + \defect{t}$. Thus $\defect{t} = 0$.
	By \cref{prelim:prop:item:defdupl2} from~\cref{prelim:prop:gentransdef},
	we have $|\dupl{t}| \le \defect{t} = 0$, i.e.~$t$ is injective.	
\end{proof} 

\begin{corollary}
    Let $t \colon X \to X$ and $S \subseteq X$ be finite.
    Then~$t$ permutes~$S$ if and only if $t(S) = S$.		
\end{corollary}

\begin{proposition} \label{prelim:prop:quasiperm}
	Let $t \colon [n] \to [n]$. Then the following are equivalent.
	\begin{enumerate}
		\item \label{prelim:prop:quasiperm:1} $t$ is a quasipermutation.	
		
		\item \label{prelim:prop:quasiperm:2} For each $x \in X$, $t(x)$ is a cycle point.
		
		\item \label{prelim:prop:quasiperm:3} For each $x \in \excl{t}$, $t(x)$ is a cycle point.	
		
		\item \label{prelim:prop:quasiperm:4} The set of tail points equals $\excl{t}$.
		
		\item \label{prelim:prop:quasiperm:5} $\dupl{t_{|t([n])}} = \emptyset$.
		
		% noch: t^{-1}(dupl) nur excl und cycle states
	\end{enumerate}	
\end{proposition}
\begin{proof}
	$\ref{prelim:prop:quasiperm:1}\Rightarrow \ref{prelim:prop:quasiperm:2}$:
	By assumption, $t_{|t(X)}$ is a permtuation.
	Then the claim follows from \cref{prelim:prop:item:cycles} of \cref{prelim:prop:permutation}
	applied to the permutation $t_{|t(X)}$.
		
	$\ref{prelim:prop:quasiperm:2}\Rightarrow \ref{prelim:prop:quasiperm:3}$:	This is clear.
		
	$\ref{prelim:prop:quasiperm:3}\Rightarrow \ref{prelim:prop:quasiperm:4}$:	
	Clearly, states in $\excl{t}$ are tail points.
	For a tail point $x$ we have $x = t^i(z)$
	for an excluded point~$z$ and some $i \ge 0$.
	By assumption, $i = 0$ and so $x = z$.
	Thus $\excl{t}$ equals precisely the tail points.	
	
	$\ref{prelim:prop:quasiperm:4}\Rightarrow \ref{prelim:prop:quasiperm:5}$:
	Assume $t_{|t(X)}$ is not injective.
	Then $t(t(x)) = t(t(z))$ for $x,z \in [n]$
	with $t(x) \ne t(z)$.
	Since a cycle point has precisely one predecessor on its cycle,
	one of both points $t(x)$ or $t(z)$ must be a tail point,
	and this point is $t(X)$.
	
	$\ref{prelim:prop:quasiperm:5}\Rightarrow \ref{prelim:prop:quasiperm:1}$:
	We have $t(t([n])) \subseteq t([n])$.
	By assumption, $t$ acts injectively on $t(X)$.
	This yields $|t(t([n]))| = |t([n])|$ and thus we 
	obtain $t(t([n])) = t([n])$.
\end{proof}

\begin{lemma} \label{prelim:trans:lem:stable-fibers}
	Let $t \colon [n] \to [n]$ be a quasipermutation. Then
	\[
	|t^{-1}(y)| = |t^{-1}(y) \cap \excl{t}| + 1
	\quad \text{for } y \in t([n]),
	\]
	and $t^{-1}(y) = \emptyset$ otherwise. In particular
	$\dupl{t} = t(\excl{t})$.
\end{lemma}
\begin{proof}
	As $t(t([n])) = t([n])$ and $t([n])$ is finite, $t_{|t([n])}$ is a permutation
	of $t([n])$. Hence every $y \in t([n])$ has exactly one preimage in $t([n])$,
	and the remaining preimages lie in $\excl{t} = [n] \setminus t([n])$, which
	gives the displayed equality; points outside $t([n])$ have no preimage.
	Consequently $|t^{-1}(y)| \ge 2$ if and only if $t^{-1}(y)$ meets $\excl{t}$,
	that is, if and only if $y \in t(\excl{t})$.
\end{proof}

\begin{proposition}
	Let $t \colon [n] \to [n]$ be of defect at most two.
	If $\dupl{t} = t(\excl{t})$, then~$t$ is a quasipermutation. 
\end{proposition}
\begin{proof}
	If there exists $x \in \dupl{t}$
	not contained in a cycle, then $x = t^i(y) = t^j(z)$
	for two distinct excluded points $y, z \in \excl{t}$ and $i,j \ge 1$.
	By \cref{prelim:prop:item:root} of \cref{prelim:prop:non-perm-dupl-cycle} 
	there exists another point in $\dupl{t}$ contained in a cycle,
	which must be the image of yet another excluded point by assumption.
	Then $t$ would have defect at least three.
	Thus every duplicate point is contained in a cycle.
	Then, with $\dupl{t} = t(\excl{t})$,
	the image of every excluded point is a cycle point.
	Thus, with \cref{prelim:prop:quasiperm}, $t$ is a quasipermutation.	
\end{proof}